%===============================================================================
% LLMWiki Technical Specification — Volume 02: Semantic Parsing
%===============================================================================
\chapter{Semantic Parsing}
\label{cha:semantic-parsing}

This volume specifies the \emph{semantic parsing} layer---the bridge between raw artifacts (source code, documents, schemas, APIs) and the content-addressed semantic units that populate the knowledge graph and search indices. Unlike traditional document chunking, semantic parsing extracts \emph{typed, structured, hierarchically organized} units with explicit semantic annotations.

\section{Design Goals}
\label{sec:parsing:goals}

\begin{designprinciple}[Language Awareness]
Parsers understand the syntax and semantics of their target language (Python, SQL, OpenAPI, Markdown, etc.), producing units that reflect logical constructs (functions, classes, tables, sections) rather than arbitrary text spans.
\end{designprinciple}

\begin{designprinciple}[Structure Preservation]
The hierarchical structure of the source artifact is preserved in the unit tree. A Python class contains method units; a Markdown document contains section units; an OpenAPI spec contains operation units.
\end{designprinciple}

\begin{designprinciple}[Content Addressability]
Each semantic unit is identified by SHA-256 of its canonical representation. Identical logical constructs yield identical identifiers regardless of source file location.
\end{designprinciple}

\begin{designprinciple}[Incremental Re-parse]
When an artifact changes, only affected units are re-parsed. The parser emits a diff of (added, modified, removed) unit IDs.
\end{designprinciple}

\begin{designprinciple}[Extensibility]
New parsers are added via the \texttt{Parser} trait (Vol.~01, \S\ref{iface:parser}) without modifying core ingestion logic.
\end{designprinciple}

\section{Parser Architecture}
\label{sec:parsing:architecture}

\begin{figure}[htbp]
\centering
\begin{adjustbox}{max width=\textwidth}
\begin{tikzpicture}[
  parser/.style={draw, rounded corners, fill=llmwikiblue!10, minimum width=3cm, minimum height=1.2cm},
  artifact/.style={draw, rounded corners, fill=llmwikigreen!10, minimum width=3cm, minimum height=1.2cm},
  unit/.style={draw, rounded corners, fill=llmwikiorange!10, minimum width=3cm, minimum height=1.2cm},
  flow/.style={->, thick, >=stealth}
]
  \node[artifact] (source) at (0,0) {Source Artifact};
  \node[parser] (detect) at (4,1.5) {MIME Detection};
  \node[parser] (route) at (4,0) {Parser Router};
  \node[parser] (parse) at (4,-1.5) {Language Parser};
  \node[unit] (units) at (8,0) {Semantic Units};
  \node[unit] (meta) at (8,2.5) {Parse Metadata};

  \draw[flow] (source) -- (detect);
  \draw[flow] (detect) -- (route);
  \draw[flow] (route) -- (parse);
  \draw[flow] (parse) -- (units);
  \draw[flow] (parse) -- (meta);
  
  % Parser registry
  \node[draw, dashed, rounded corners, align=left, font=\small] (registry) at (10,-1.5) {%
    \textbf{Parser Registry}\\
    python : TreeSitterParser\\
    markdown : MarkdownParser\\
    pdf : PdfParser\\
    openapi : SpecParser\\
    sql : SqlParser\\
    yaml/json : StructuredParser\\
    docx/pptx/xlsx : OoxmlParser\\
    hwpx : HwpxParser\\
    hwp : HwpParser (read-only)};
  \draw[flow, dashed] (route) -- (registry);
\end{tikzpicture}
\end{adjustbox}
\caption{Parser pipeline: detection, routing, parsing, and unit emission.}
\label{fig:parsing:pipeline}
\end{figure}

\subsection{Parser Trait}
\label{sec:parsing:parser-trait}

Every parser implements the \texttt{Parser} interface (Vol.~01, \S\ref{iface:parser}):

\begin{verbatim}
trait Parser {
    fn parse(&self, artifact: Artifact) -> Result<Vec<SemanticUnit>>;
    fn supported_mimes(&self) -> Vec<MimeType>;
    fn schema(&self) -> ParserSchema;
}
\end{verbatim}

\textbf{Artifact} encapsulates the raw bytes, path, MIME type, and optional metadata (encoding, language hints).
\textbf{SemanticUnit} is the canonical output (defined below).
\textbf{ParserSchema} describes the unit types this parser produces (e.g., \texttt{PythonFunction}, \texttt{SQLTable}, \texttt{MarkdownSection}).

\subsection{Parser Registry}
\label{sec:parsing:registry}

The registry maps MIME types (and file extensions as fallback) to parser instances. Parsers are loaded dynamically at startup from \texttt{config.toml} \texttt{[ingestion.parsers]}.

\begin{verbatim}
# Config example
[ingestion.parsers]
python = { type = "tree-sitter", grammars = ["python"] }
markdown = { type = "markdown", extract_frontmatter = true }
pdf = { type = "pdf", ocr = false }
openapi = { type = "openapi", version = "3.1" }
sql = { type = "sql", dialect = "postgres" }
docx = { type = "ooxml", kind = "wordprocessing", track_changes = "as-metadata" }
pptx = { type = "ooxml", kind = "presentation" }
xlsx = { type = "ooxml", kind = "spreadsheet" }
hwpx = { type = "hwpx" }
hwp  = { type = "hwp", mode = "read-only" }  # see Sec. parsing:office-formats
\end{verbatim}

\subsection{Office and HWP Format Coverage}
\label{sec:parsing:office-formats}

DOCX, PPTX, XLSX, and HWPX are all XML-in-a-zip container formats (OOXML for the first three, the HWPX schema for the fourth); this section states what \texttt{OoxmlParser}/\texttt{HwpxParser} extract as \texttt{UnitOccurrence}s and how the corresponding \texttt{DocumentWriter} (Vol.~12 \S\ref{sec:mutation:interfaces}) writes back, since read and write coverage are tracked separately (Vol.~12 \S\ref{sec:mutation:formats} is authoritative on write support; this section is read-side).

\begin{itemize}
  \item \textbf{DOCX}: paragraphs and headings as \texttt{Section} units; tables as \texttt{Table} units with cell-level \texttt{UnitOccurrence}s; embedded comments and tracked-changes (\texttt{w:ins}/\texttt{w:del}) are extracted as metadata attached to the paragraph they annotate, not discarded --- an evidence-gate citation into a document under active review should be able to say whether the cited text is an accepted change, a pending insertion, or a pending deletion.
  \item \textbf{PPTX}: each slide is a unit; text boxes, tables, and speaker notes within a slide are child units (span type \texttt{slide\_object\_id}, Appendix~B \texttt{UnitOccurrence}); embedded charts/images are recorded as opaque objects with position metadata (\S\ref{sec:mutation:formats}), not parsed for content.
  \item \textbf{XLSX}: each sheet is a unit; cells and named ranges are child units (span type \texttt{sheet\_cell}); formulas are extracted as both their formula text and (when the workbook carries a cached value) their last-computed value, tagged separately so retrieval can distinguish ``the formula is'' from ``the value was.''
  \item \textbf{HWPX}: the XML-based Hangul Word Processor format; structurally closer to DOCX (paragraph/section-based) and parsed on the same principles above via \texttt{HwpxParser}.
  \item \textbf{HWP (binary)}: read-only in this revision. The legacy binary HWP format has no publicly specified write-back path that preserves round-trip fidelity comparable to OOXML/HWPX; \texttt{HwpParser} extracts \texttt{UnitOccurrence}s for retrieval, but Vol.~12's approval pipeline does not offer HWP as a mutation target until a \texttt{DocumentWriter} for it exists and passes \texttt{FormatValidator} (Vol.~12 \S\ref{sec:mutation:formats}). This is a stated gap, not a silent one: a mutation request against a \texttt{.hwp} file is rejected at \texttt{PatchPlanner} with an explicit ``format not writable'' error rather than silently falling back to a lossy conversion.
\end{itemize}

\section{Semantic Unit Model}
\label{sec:parsing:semantic-unit}

\begin{definition}[Semantic Unit and Its Occurrence]
\label{def:semantic-unit}
A parser emits two related but distinct records per parsed fragment, not one (Appendix~B \S\ref{app:data-model:identity} gives the full schema; this is the parser-facing view of the same split):
\begin{equation}
u = (\text{id}, \text{type}, \text{content}, \text{metadata}, \text{children}), \qquad
o = (\text{occurrence\_id}, \text{content\_id}, \text{span}, \text{artifact\_version\_id})
\end{equation}
where:
\begin{itemize}
  \item \texttt{id} (= \texttt{content\_id}) = SHA-256(canonical(content, type)) --- \textbf{not} a function of \texttt{span}. Two occurrences of identical content at different spans (the same helper function copy-pasted into two files, a repeated boilerplate header) share this id; that sharing is what makes deduplication work, and it is destroyed if span is hashed into the id.
  \item \texttt{type} $\in$ \{\texttt{Function}, \texttt{Class}, \texttt{Section}, \texttt{Table}, \texttt{Operation}, \dots\}
  \item \texttt{content}: canonical text representation (whitespace-normalized)
  \item \texttt{metadata}: typed key-value annotations (see \S\ref{sec:parsing:metadata})
  \item \texttt{children}: ordered list of child unit IDs (hierarchy)
  \item \texttt{occurrence\_id} = SHA-256(artifact\_version\_id, span) --- unique per location, even when \texttt{content\_id} repeats
  \item \texttt{span}: \texttt{(start\_line, end\_line, start\_col, end\_col)} in source for text/code; other formats use other span types (Appendix~B \texttt{UnitOccurrence.span\_type})
\end{itemize}
A parser that only emits $u$ and discards $o$ (or that folds span into \texttt{id} as an earlier draft of this definition did) cannot answer ``where did this content come from'' once a second occurrence exists, which is exactly the failure mode Appendix~B \S\ref{app:data-model:identity} documents.
\end{definition}

\subsection{Unit Type Taxonomy}
\label{sec:parsing:unit-types}

\begin{table}[htbp]
\centering
\caption{Core Semantic Unit Types by Parser Domain}
\label{tab:parsing:unit-types}
\begin{adjustbox}{max width=\textwidth}
\begin{tabular}{lll}
\toprule
\textbf{Domain} & \textbf{Unit Types} & \textbf{Key Metadata} \\
\midrule
Python & Module, Class, Function, Method, Decorator, Import & signature, decorators, type\_hints, docstring \\
JavaScript/TS & Module, Class, Function, ArrowFunction, Interface, Type & signature, jsdoc, exports \\
Markdown & Document, Section, CodeBlock, List, Table, Link & heading\_level, language, alt\_text \\
PDF & Page, Paragraph, Table, Figure, Header, Footer & bbox, font\_info, page\_num \\
OpenAPI & Info, Server, Path, Operation, Schema, Parameter & method, path, status\_codes, auth \\
SQL DDL & Database, Schema, Table, Column, Index, Constraint & data\_type, nullable, default, references \\
YAML/JSON & Document, Mapping, Sequence, Scalar & json\_path, schema\_ref \\
Dockerfile & Instruction, Stage, Arg, Env, Label & instruction, value, stage \\
Protobuf & Message, Enum, Service, Field, RPC & field\_number, type, label \\
\bottomrule
\end{tabular}
\end{adjustbox}
\end{table}

\subsection{Metadata Schema}
\label{sec:parsing:metadata}

Each unit type defines a typed metadata schema. Examples:

\begin{verbatim}
# Python Function metadata
{
  "signature": "def foo(x: int, y: str = 'a') -> bool",
  "decorators": ["@staticmethod", "@cache"],
  "type_hints": {"x": "int", "y": "str", "return": "bool"},
  "docstring": "Returns True if...",
  "is_async": false,
  "is_generator": false,
  "complexity": 3
}

# SQL Table metadata
{
  "columns": [
    {"name": "id", "type": "BIGINT", "nullable": false, "primary_key": true},
    {"name": "email", "type": "VARCHAR(255)", "nullable": false, "unique": true}
  ],
  "indexes": [{"name": "idx_email", "columns": ["email"], "unique": true}],
  "foreign_keys": [{"column": "user_id", "references": "users.id"}],
  "engine": "InnoDB",
  "charset": "utf8mb4"
}

# OpenAPI Operation metadata
{
  "method": "POST",
  "path": "/api/v1/users",
  "operation_id": "createUser",
  "parameters": [{"name": "body", "in": "body", "schema": "UserInput"}],
  "responses": {"201": "User", "400": "Error"},
  "security": [{"BearerAuth": []}],
  "deprecated": false
}
\end{verbatim}

\subsection{Canonicalization}
\label{sec:parsing:canonicalize}

Content-addressability requires deterministic canonicalization:

\begin{algorithm}[H]
\caption{Unit Canonicalization}
\label{alg:canonicalize}
\begin{algorithmic}[1]
\Require Unit $u$ with fields $(type, content, span, metadata, children)$
\Ensure Canonical byte string $C$ for hashing
\State $C \gets \text{type} \parallel \texttt{:}$ \Comment{Type prefix}
\State $C \gets C \parallel \textsc{NormalizeWhitespace}(content)$ \Comment{Canonical text}
\State $C \gets C \parallel \textsc{SortKeys}(metadata)$ \Comment{Deterministic JSON}
\State $C \gets C \parallel \textsc{Sort}(children)$ \Comment{Sorted child IDs}
\State \Return $C$
\end{algorithmic}
\end{algorithm}

This ensures: (a) semantically identical units in different files produce identical IDs, (b) formatting changes that don't affect semantics (whitespace, comment reordering) don't change IDs if metadata excludes them.

\section{Language-Specific Parsers}
\label{sec:parsing:language-parsers}

\subsection{Tree-sitter Based Parsers (Code)}
\label{sec:parsing:tree-sitter}

For programming languages, we use \href{https://tree-sitter.github.io/}{tree-sitter} for robust, incremental, error-tolerant parsing.

\begin{designdecision}
\label{dec:tree-sitter}
\textbf{Decision:} Use tree-sitter as the primary parsing engine for all supported programming languages.
\end{designdecision}

\begin{designrationale}
\begin{enumerate}
  \item \textbf{Incremental re-parse}: tree-sitter reuses syntax tree nodes for unchanged regions ($O(\log n)$ edits).
  \item \textbf{Error tolerance}: Recovers gracefully from syntax errors; produces partial trees.
  \item \textbf{Unified API}: Same query language (\texttt{SCM}) across 40+ languages.
  \item \textbf{Bindings}: Rust (\texttt{tree-sitter}), Python (\texttt{tree-sitter-languages}), Node.js available.
  \item \textbf{Grammar maintenance}: Community-maintained grammars; we vendor specific versions.
\end{enumerate}
\end{designrationale}

\paragraph{Tree-sitter Query Patterns.}
Semantic units are extracted via SCM queries. Example for Python functions:

\begin{verbatim}
(function_definition
  name: (identifier) @name
  parameters: (parameters) @params
  return_type: (type)? @return_type
  body: (block) @body
  decorators: (decorator)* @decorators
) @function
\end{verbatim}

The query captures are mapped to \texttt{SemanticUnit} fields. A \texttt{QueryMapper} per language handles this transformation.

\subsubsection{Supported Languages}
\label{sec:parsing:ts-languages}

\begin{table}[htbp]
\centering
\caption{Tree-sitter Languages Supported (v0.1)}
\label{tab:parsing:ts-languages}
\begin{adjustbox}{max width=\textwidth}
\begin{tabular}{lll}
\toprule
\textbf{Language} & \textbf{Grammar Version} & \textbf{Unit Types} \\
\midrule
Python & 0.20.1 & Module, Class, Function, Method, Decorator, Import, TypeAlias \\
JavaScript & 0.20.0 & Module, Class, Function, ArrowFunction, Interface, Type, Import \\
TypeScript & 0.20.0 & (extends JS) + TypeAnnotation, Generic \\
Rust & 0.20.1 & Module, Struct, Enum, Fn, Impl, Trait, Macro, Use \\
Go & 0.20.0 & Package, Struct, Func, Method, Interface, Type, Import \\
Java & 0.20.1 & Class, Interface, Method, Field, Constructor, Annotation \\
C/C++ & 0.20.1 & Function, Struct, Class, Macro, Include, Namespace \\
C\# & 0.20.0 & Class, Method, Property, Interface, Attribute, Using \\
SQL & 0.20.0 & (see \S\ref{sec:parsing:sql-parser}) \\
\bottomrule
\end{tabular}
\end{adjustbox}
\end{table}

\subsection{Markdown Parser}
\label{sec:parsing:markdown}

Uses \texttt{pulldown-cmark} (Rust) or \texttt{markdown-it} (Python) with extensions:
\begin{itemize}
  \item GitHub Flavored Markdown (GFM): tables, task lists, strikethrough
  \item Frontmatter (YAML/TOML/JSON) extraction
  \item Heading anchoring and TOC generation
  \item Code block language detection
  \item Link/reference resolution
\end{itemize}

Units produced: \texttt{Document}, \texttt{Section}(level), \texttt{Paragraph}, \texttt{CodeBlock}(lang), \texttt{Table}, \texttt{List}, \texttt{Link}, \texttt{Image}.

\subsection{PDF Parser}
\label{sec:parsing:pdf}

Uses \texttt{pdfplumber} (Python) or \texttt{lopdf} (Rust) with optional OCR (\texttt{tesseract}):
\begin{itemize}
  \item Text extraction with position (bbox, font, size)
  \item Table detection (via \texttt{pdfplumber} table extraction)
  \item Figure/image extraction (bbox, optional OCR)
  \item Reading order reconstruction (column-aware)
\end{itemize}

Units: \texttt{Page}, \texttt{Paragraph}, \texttt{Table}, \texttt{Figure}, \texttt{Header}, \texttt{Footer}.

\subsection{OpenAPI / AsyncAPI Parser}
\label{sec:parsing:openapi}

Parses OpenAPI 3.0/3.1 and AsyncAPI 2.x/3.x specs. Produces units for every reusable component:
\begin{itemize}
  \item \texttt{Info}, \texttt{Server}, \texttt{Path}, \texttt{Operation}
  \item \texttt{Schema} (with JSON Schema resolution)
  \item \texttt{Parameter}, \texttt{RequestBody}, \texttt{Response}
  \item \texttt{SecurityScheme}, \texttt{Tag}, \texttt{ExternalDocs}
\end{itemize}

Cross-references (\texttt{\$ref}) are resolved and stored as graph edges in the semantic units.

\subsection{SQL DDL Parser}
\label{sec:parsing:sql-parser}

Uses \texttt{sqlglot} (Python) for multi-dialect parsing (PostgreSQL, MySQL, SQLite, BigQuery, Snowflake, T-SQL).
Produces units for: \texttt{Database}, \texttt{Schema}, \texttt{Table}, \texttt{Column}, \texttt{Index}, \texttt{Constraint}, \texttt{View}, \texttt{Function}, \texttt{Trigger}.

Column-level lineage is extracted: \texttt{Column} units reference source columns in \texttt{SELECT} clauses.

\subsection{Structured Data Parser (YAML/JSON/TOML)}
\label{sec:parsing:structured}

For configuration and data files:
\begin{itemize}
  \item JSON Schema / OpenAPI schema detection (if \texttt{\$schema} present)
  \item Path-based units: each mapping key becomes a unit with JSON path metadata
  \item Array elements indexed with position
  \item Merge with schema annotations when available
\end{itemize}

\section{Structure-Aware Chunking}
\label{sec:parsing:chunking}

After parsing, semantic units may exceed embedding model context windows. The \texttt{Chunker} splits units while preserving semantic boundaries.

\begin{interface}[Chunker]
\label{iface:chunker-detail}
\begin{verbatim}
trait Chunker {
    fn chunk(&self, unit: SemanticUnit, config: ChunkConfig) -> Vec<Chunk>;
    fn config_schema(&self) -> ChunkConfigSchema;
}

struct ChunkConfig {
    max_tokens: usize,
    overlap_tokens: usize,
    respect_boundaries: bool,      // Don't split inside function/class
    preserve_hierarchy: bool,      // Include parent context
    strategy: ChunkStrategy,       // "hierarchical", "sliding", "semantic"
}
\end{verbatim}
\end{interface}

\subsection{Hierarchical Chunking Algorithm}
\label{sec:parsing:hierarchical-chunk}

Default strategy for code and structured documents:

\begin{algorithm}[H]
\caption{Hierarchical Chunking}
\label{alg:hierarchical-chunk}
\begin{algorithmic}[1]
\Require Unit $u$ with children $c_1, \dots, c_n$, config $cfg$
\Ensure Chunks $\mathcal{C}$ each $\leq cfg.max\_tokens$
\State $\mathcal{C} \gets \varnothing$
\State $current \gets \varnothing$, $current\_tokens \gets 0$
\ForEach{child $c \in children(u)$ in order}
  \State $child\_tokens \gets \textsc{TokenCount}(c.content)$
  \If{$current\_tokens + child\_tokens > cfg.max\_tokens$ and $current \neq \varnothing$}
    \State $\mathcal{C} \gets \mathcal{C} \cup \{current\}$
    \If{$cfg.overlap\_tokens > 0$}
      \State $current \gets \textsc{Tail}(current, cfg.overlap\_tokens)$
      \State $current\_tokens \gets \textsc{TokenCount}(current)$
    \Else
      \State $current \gets \varnothing$, $current\_tokens \gets 0$
    \EndIf
  \EndIf
  \State $current \gets current \parallel c.content$
  \State $current\_tokens \gets current\_tokens + child\_tokens$
\EndFor
\If{$current \neq \varnothing$}
  \State $\mathcal{C} \gets \mathcal{C} \cup \{current\}$
\EndIf
\State \Return $\mathcal{C}$
\end{algorithmic}
\end{algorithm}

\subsection{Sliding Window (Fallback)}
\label{sec:parsing:sliding}

For plain text or when hierarchical structure is unavailable:

\begin{verbatim}
# Sliding window with sentence boundary awareness
strategy = "sliding"
window_size = 512      # tokens
overlap = 50           # tokens
split_on = ["\n\n", "\n", ". ", "! ", "? ", "; "]
\end{verbatim}

\subsection{Semantic Chunking (Experimental)}
\label{sec:parsing:semantic-chunk}

Uses embedding similarity to find natural boundaries:

\begin{algorithm}[H]
\caption{Semantic Chunking via Embedding Similarity}
\label{alg:semantic-chunk}
\begin{algorithmic}[1]
\Require Sentences $s_1, \dots, s_n$, threshold $\theta$, max\_tokens
\Ensure Chunks $\mathcal{C}$
\State $embs \gets \textsc{Embed}(s_1, \dots, s_n)$
\State $boundaries \gets \{1\}$
\For{$i = 2$ to $n$}
  \State $sim \gets \textsc{Cosine}(embs[i-1], embs[i])$
  \If{$sim < \theta$}
    \State $boundaries \gets boundaries \cup \{i\}$
  \EndIf
\EndFor
\State $boundaries \gets boundaries \cup \{n+1\}$
\State $\mathcal{C} \gets \varnothing$
\ForEach{consecutive $(b_j, b_{j+1}) \in boundaries$}
  \State $chunk \gets \textsc{MergeSentences}(s_{b_j}, \dots, s_{b_{j+1}-1})$
  \If{$\textsc{TokenCount}(chunk) > max\_tokens$}
    \State $\mathcal{C} \gets \mathcal{C} \cup \textsc{SlidingWindow}(chunk)$
  \Else
    \State $\mathcal{C} \gets \mathcal{C} \cup \{chunk\}$
  \EndIf
\EndFor
\State \Return $\mathcal{C}$
\end{algorithmic}
\end{algorithm}

\subsection{Scale-Free Chunking by Graph Renormalization}
\label{sec:parsing:renormalization}

The three strategies above share a defect that no amount of tuning removes: each commits to a scale. \texttt{window\_size = 512} is a claim that the natural unit of meaning in the corpus is 512 tokens wide, and that claim is false for every corpus --- a call site, a function, a module, and a subsystem are all legitimate retrieval targets at different query granularities. Fixing the scale at ingestion time forces the retrieval stage to reconstruct the levels that chunking destroyed.

Coarse-graining across scales without committing to one is the subject of the renormalization group. The construction repeatedly integrates out short-range degrees of freedom, generating a flow through a family of effective descriptions in which structures that survive many steps are the physically meaningful ones. Recent work carries this to graphs directly: \citet{spectralcoarse2025} define a renormalization procedure over the coarse-grained Laplacian that produces reduced representations across scales while preserving both dynamics and large-scale topology, and note explicitly that this differs from clustering methods such as Louvain --- renormalization captures characteristic structure \emph{across} scales rather than partitioning at one.

Since LLMWiki already maintains a knowledge graph over units (Vol.~03), the machinery applies without a new data structure.

\begin{definition}[Renormalization Chunk Ladder]
\label{def:chunk-ladder}
Let $G_0$ be the unit graph with edges weighted by containment, adjacency, and reference. A chunk ladder is the sequence $G_0 \to G_1 \to \dots \to G_L$ where each $G_{\ell+1}$ is obtained by spectral coarse-graining of $G_{\ell}$, and the chunk set at level $\ell$ is the set of vertices of $G_{\ell}$, each carrying the concatenated content of the units it subsumes.
\end{definition}

\begin{observation}[Persistence Selects Boundaries]
\label{obs:persistence}
A chunk boundary that survives many renormalization steps separates parts of the corpus that remain weakly coupled under coarse-graining; a boundary that dissolves after one step was an artifact of local token statistics. Boundary quality is therefore measurable as persistence across the ladder, replacing the unfalsifiable question ``is 512 the right window?'' with the answerable one ``at which level does this boundary disappear?''
\end{observation}

Retrieval indexes every level of the ladder, and the query planner selects the level rather than the ingestion configuration doing so: a query naming a specific symbol resolves against $G_0$, while a query about a subsystem's responsibilities resolves against a coarse level whose vertices correspond to that subsystem. The cost is index size, which grows as $\sum_{\ell} |V(G_\ell)|$ --- bounded by a small multiple of $|V(G_0)|$ because coarse-graining contracts the vertex set geometrically.

\paragraph{Status.} This is the least mature of the mechanisms specified in this document. Spectral coarse-graining of the unit graph is well-defined and implementable, but the claim that renormalization-persistent boundaries yield better retrieval than tuned fixed-width chunking is an empirical one that the evaluation framework of Vol.~06 must settle before it is relied upon. It is specified here because the ladder subsumes the existing strategies as special cases --- $L=0$ recovers hierarchical chunking exactly --- and can therefore be adopted without discarding the current pipeline.

\section{Entity and Relation Extraction}
\label{sec:parsing:entity-extraction}

After chunking, lightweight extraction identifies entities and relations within and across units.

\subsection{Rule-Based Extraction}
\label{sec:parsing:rule-extraction}

Zero-LLM approach using tree-sitter queries and regex patterns:
\begin{itemize}
  \item \textbf{Code}: Function calls, imports, class inheritance, type references
  \item \textbf{Markdown}: Links, cross-references, image references
  \item \textbf{OpenAPI}: Schema references, parameter bindings, auth requirements
  \item \textbf{SQL}: Table/column references, foreign keys, join conditions
  \item \textbf{Config}: Environment variable references, file path references
\end{itemize}

Produces \texttt{(source\_unit, target\_unit, relation\_type, confidence)} tuples.

\subsection{LLM-Assisted Extraction (Optional)}
\label{sec:parsing:llm-extraction}

For complex relations (e.g., implicit semantic links, cross-file call graphs), a local LLM can be invoked:

\begin{verbatim}
# Extraction prompt template (Vol. 05 Agent Runtime)
EXTRACT_ENTITIES_RELATIONS = """
Given the following semantic units from a {language} codebase,
extract entities and relations in JSON format.

Units:
{units_json}

Entity types: {config.entity_types}
Relation types: {config.relation_types}

Output format:
{{
  "entities": [{"id": "...", "type": "...", "name": "...", "unit_id": "..."}],
  "relations": [{"source": "...", "target": "...", "type": "...", "confidence": 0.0-1.0}]
}}
"""
\end{verbatim}

\subsection{Entity Resolution}
\label{sec:parsing:entity-resolution}

Extracted entities are resolved against the knowledge graph (Vol.~03):

\begin{algorithm}[H]
\caption{Entity Resolution}
\label{alg:entity-resolution-parsing}
\begin{algorithmic}[1]
\Require Extracted entity $e$ with name $n$, type $t$; graph $G$
\Ensure Resolved node ID or new node
\State $candidates \gets \textsc{GraphLookup}(G, name=n, type=t)$
\If{$candidates = \varnothing$}
  \State \Return \textsc{CreateNode}(G, e) \Comment{New entity}
\Else
  \State $best \gets \arg\max_{c \in candidates} \textsc{Similarity}(c, e)$
  \If{$\textsc{Similarity}(best, e) \geq \tau_{er}$}
    \State \Return $best.id$ \Comment{Merge}
  \Else
    \State \Return \textsc{CreateNode}(G, e) \Comment{Distinct entity}
  \EndIf
\EndIf
\end{algorithmic}
\end{algorithm}

Similarity combines: name fuzzy match, type compatibility, context overlap (shared neighbors), embedding cosine similarity.

\section{Incremental Parsing}
\label{sec:parsing:incremental}

When a file changes, re-parsing only affected regions is critical for latency.

\paragraph{Tree-sitter Incremental Parse.}
Tree-sitter accepts the old tree and the edit (byte range + new text). It returns a new tree with reused nodes. We compute the diff at unit level:

\begin{algorithm}[H]
\caption{Incremental Unit Diff}
\label{alg:incremental-diff}
\begin{algorithmic}[1]
\Require Old units $U_{old}$, new tree $T_{new}$, edit range $E$
\Ensure $(added, modified, removed)$ unit ID sets
\State $U_{new} \gets \textsc{ExtractUnits}(T_{new})$
\State $affected \gets \{ u \in U_{old} \mid u.span \cap E \neq \varnothing \}$
\State $unchanged \gets U_{old} \setminus affected$
\State $modified \gets \{ u_{new} \in U_{new} \mid \exists u_{old} \in affected: u_{new}.id = u_{old}.id \land u_{new} \neq u_{old} \}$
\State $added \gets \{ u_{new} \in U_{new} \mid u_{new}.id \notin U_{old}.ids \}$
\State $removed \gets \{ u_{old} \in affected \mid u_{old}.id \notin U_{new}.ids \}$
\State \Return $(added, modified, removed)$
\end{algorithmic}
\end{algorithm}

Only $added \cup modified$ units are re-embedded and re-indexed. $removed$ units are tombstoned in indices.

\section{Embedding Generation}
\label{sec:parsing:embedding}

The \texttt{Embedder} trait (Vol.~01, \S\ref{iface:embedder}) abstracts local embedding models.

\subsection{Supported Models}
\label{sec:parsing:embed-models}

\begin{table}[htbp]
\centering
\caption{Local Embedding Models (Recommended)}
\label{tab:parsing:embed-models}
\begin{adjustbox}{max width=\textwidth}
\begin{tabular}{llcccc}
\toprule
\textbf{Model} & \textbf{Dim} & \textbf{Max Tok} & \textbf{Speed (CPU)} & \textbf{MTEB} & \textbf{License} \\
\midrule
BGE-M3 & 1024 & 8192 & $\sim$500 tok/s & 62.3 & MIT \\
E5-Large-v2 & 1024 & 512 & $\sim$400 tok/s & 61.5 & MIT \\
Nomic-Embed-Text-v1.5 & 768 & 8192 & $\sim$600 tok/s & 60.1 & Apache-2.0 \\
BGE-Small-En-v1.5 & 384 & 512 & $\sim$1200 tok/s & 58.2 & MIT \\
Multilingual-E5-Large & 1024 & 512 & $\sim$350 tok/s & 63.1 & MIT \\
\bottomrule
\end{tabular}
\end{adjustbox}
\end{table}

\subsection{Batch Inference}
\label{sec:parsing:batch-inference}

All embeddings use batched inference for throughput:

\begin{verbatim}
# Embedder config
batch_size = 32
max_seq_len = 512
normalize = true
device = "cpu"  # or "cuda", "mps"
onnx_runtime = true  # optimized inference
\end{verbatim}

Embeddings are stored as \texttt{float32} in the HNSW index (Vol.~03) and as \texttt{float16} in the unit metadata for memory efficiency.

\section{Parser Extensibility Guide}
\label{sec:parsing:extensibility}

To add a custom parser:

1. Implement \texttt{Parser} trait
2. Register in \texttt{config.toml}
3. Define unit types and metadata schemas
4. (Optional) Add tree-sitter grammar or SCM queries
5. Write tests against \texttt{testdata/} fixtures

\begin{verbatim}
# Example: Custom parser for GraphQL SDL
[ingestion.parsers.graphql]
type = "custom"
module = "my_parsers.graphql_parser"
class = "GraphQLParser"
supported_mimes = ["application/graphql", "text/x.graphql"]
extensions = [".graphql", ".gql"]
\end{verbatim}

\section{Summary}
\label{sec:parsing:summary}

The semantic parsing layer transforms heterogeneous artifacts into a uniform representation of content-addressed, hierarchically structured, metadata-rich semantic units. Key properties:

\begin{enumerate}
  \item \textbf{Language-aware}: Tree-sitter for code, specialized parsers for structured formats
  \item \textbf{Hierarchical}: Units form trees reflecting source structure
  \item \textbf{Content-addressed}: SHA-256 identifiers enable deduplication and lineage
  \item \textbf{Incremental}: Tree-sitter's incremental parsing + unit-level diff
  \item \textbf{Extensible}: Parser trait + registry for new formats
  \item \textbf{Metadata-rich}: Typed annotations per unit type for downstream use
\end{enumerate}

The next volume (Vol.~03) details how these units populate the knowledge graph and semantic indices.

%===============================================================================
% End of Volume 02
%===============================================================================
